\documentclass{article}
\usepackage[UTF8]{ctex}
\usepackage{amsmath,amssymb,amsfonts}
\usepackage{enumitem}
\usepackage{geometry}
\geometry{a4paper, margin=1in}
% === 错题本排版优化（不参与题目标识） ===
\setlength{\parindent}{0pt}
\setlength{\parskip}{0.35em plus 0.10em minus 0.08em}
\setlength{\emergencystretch}{3em}
\setlist{itemsep=0.25em, topsep=0.40em, parsep=0pt}
\allowdisplaybreaks[2]
% === 错题本排版优化结束 ===
\begin{document}
\section*{长安大学2024年高等代数真题}
\section*{174．长安大学2024年高等代数试卷}
一．填空题．每题5分，共30分．\\

\subsection{题目1}
1．设 \((x-1)^{2} \mid A x^{4}+B x^{2}+1\) ，则 \(A=\) \(\_\_\_\_\)，\(B=\) \(\_\_\_\_\)．

\subsection{题目2}
2．\(n\) 阶行列式 \(D=\left|d_{i j}\right|\) ，其中 \(d_{i j}=\left\{\begin{array}{cc}-1, & i>j ; \\ 1, & i \leq j \text { 。则 } D \text { 的展开式中正项总数为 }\end{array}\right.\) \(\_\_\_\_\)．

\subsection{题目3}
3．设 \(A\) 为3阶方阵，若 \(A+E, A+2 E, A+3 E\) 均为退化矩阵，则 \(|2 A+8 E|=\) \(\_\_\_\_\)．\\
\(\_\_\_\_\)

\subsection{题目4}
4．当 \(a=\) \(\_\_\_\_\)时，二次型 \(f(x, y, z)=a\left(x^{2}+y^{2}+z^{2}\right)+2 x y+2 x z-2 y z\) 是负定的．

\subsection{题目5}
5．设 \(A, B, A+B\) 均为 \(n\) 阶可逆方阵，则 \(A^{-1}+B^{-1}\) 的逆矩阵为 \(\_\_\_\_\)．

\subsection{题目6}
6．设 \(A=\left(\begin{array}{ccc}1 & -1 & 1 \\ 2 & 4 & -2 \\ -3 & -3 & a\end{array}\right), B=\left(\begin{array}{lll}2 & 0 & 0 \\ 0 & 2 & 0 \\ 0 & 0 & b\end{array}\right)\) ，则当 \(P=\) \(\_\_\_\_\)时，满足 \(P^{-1} A P=B\) 。

二．（10 分）应用一元多项式理论证明 \(\sin 10^{\circ}\) 为无理数．

三．（10 分）设 \(x_{1} x_{2} \cdots x_{n} \neq 0\) ，计算下面 \(n\) 阶行列式

\[
D=\left|\begin{array}{ccccc}
1 & 1 & 1 & \cdots & 1 \\
x_{1}^{2} & x_{2}^{2} & x_{3}^{2} & \cdots & x_{n}^{2} \\
x_{1}^{3} & x_{2}^{3} & x_{3}^{3} & \cdots & x_{n}^{3} \\
\vdots & \vdots & \vdots & & \vdots \\
x_{1}^{n} & x_{2}^{n} & x_{3}^{n} & \cdots & x_{n}^{n}
\end{array}\right| .
\]

四．（20 分）设 \(c \neq 0\) ，求下面线性方程组的通解．

\[
\begin{cases}x_{1}+x_{2} & =c_{1} \\ x_{2}+x_{3} & =c_{2} \\ \cdots \cdots & \\ x_{n-1}+x_{n} & =c_{n} \\ x_{1}+x_{2} & =c_{n}\end{cases}
\]

五．（10分）设 \(A, B\) 是 \(n\) 阶方阵，且满足 \(A B=A+B\) ，证明：\(A B=B A\) 。

六．（20 分）已知矩阵

\[
A=\left(\begin{array}{cccc}
1 & 1 & \cdots & 1 \\
1 & 1 & \cdots & 1 \\
\vdots & \vdots & & \vdots \\
1 & 1 & \cdots & 1
\end{array}\right)_{n \times n} .
\]

\begin{itemize}
  \item[（1）] 求满足 \(f(A)=O\) 的一个二次多项式 \(f(x)\) ．
  \item[（2）] 求 \(A^{100}\) 。
  \item[（3）] 求 \(\left(E_{n}+A\right)^{100}\) ．
  \item[（4）] 求 \(\left(E_{n}+A\right)^{-1}\) ．
\end{itemize}

七．（20 分）设 \(\mathcal{A}\) 是数域 \(P\) 上 \(n\) 维线性空间 \(V\) 的一个线性变换．证明：

\begin{itemize}
  \item[（1）] 在 \(P[x]\) 中有一个次数不超过 \(n^{2}\) 的多项式 \(f(x)\) ，使得 \(f(\mathcal{A})=\mathcal{O}\) ．
  \item[（2）] 如果 \(f(\mathcal{A})=\mathcal{O}, g(\mathcal{A})=\mathcal{O}\) ，那么 \(d(\mathcal{A})=\mathcal{O}\) ，这里 \(d(x)\) 是 \(f(x)\) 与 \(g(x)\) 的最大公因式．
  \item[（3）] \(\mathcal{A}\) 可逆的充要条件是存在一个常数项不为 0 的多项式 \(f(x)\) 使得 \(f(\mathcal{A})=\mathcal{O}\) ．
\end{itemize}

八．（15 分）设 \(A=\left(\begin{array}{ccc}-1 & 1 & 0 \\ -4 & 3 & 0 \\ 1 & 0 & 2\end{array}\right), B=\left(\begin{array}{ccc}1 & -1 & 0 \\ 1 & 2 & 1 \\ 0 & 1 & 1\end{array}\right)\) ，证明：\(A\) 与 \(B\) 在复数域上相似．

九．（15 分）已知二次型 \(f\left(x_{1}, x_{2}, x_{3}\right)=(1-a) x_{1}^{2}+(1-a) x_{2}^{2}+2 x_{3}^{2}+2(1+a) x_{1} x_{2}\) 的秩为 2 ．

\begin{itemize}
  \item[（1）] 求 \(a\) 的值。
  \item[（2）] 用正交变换化 \(f\left(x_{1}, x_{2}, x_{3}\right)\) 为标准形（要求写出所用的正交变换）。
  \item[（3）] 求方程 \(f\left(x_{1}, x_{2}, x_{3}\right)=0\) 的解集合．
\end{itemize}
\end{document}
